\chapter{الأعضاء في العمل}\label{ch:g7:organs-at-work}

اصعد ثلاث طبقات من الدرج عدوًا ثم تفقّد حالك: \omterm{def:g1:my-body:parts}{رجلان} تحترقان، \omterm{def:g4:heart-and-blood:heart}{وقلب}
يدقّ، ونفَس يتسارع --- \omterm{def:g1:the-five-senses:sense}{وجلد}، على غرابة ذلك، محمرّ دافئ. لم يأمر
أحد بهذه التغيرات؛ بل أمرت نفسها، وبتناسب تامّ. وهذا الفصل يقيس ما
يستهلكه العضو العامل حقًّا، وكيف يعيد الجسم تنظيم توصيلاته لحظة
بدء العمل.

\section{ماذا تستهلك العضلة العاملة}

\begin{proposition}[فاتورة العضلة]\label{prop:g7:organs-at-work:bill}
\omterm{def:g3:bones-and-muscles:muscle}{العضلة} العاملة تنفق، ونفقتها قابلة للقياس. فبالقياس إلى \omterm{def:g3:bones-and-muscles:muscle}{العضلة}
نفسها مستريحةً، تكون \omterm{def:g3:bones-and-muscles:muscle}{العضلة} الشديدة العمل:
\begin{itemize}
  \item تأخذ من \emph{\omterm{def:g7:what-is-respiration:respiration}{الأكسجين}} في \omterm{prop:g4:journey-of-food:absorb}{الدم} العابر بها أكثر بكثير ---
        أضعاف سحبها في الراحة؛
  \item وتأخذ من \emph{الغلوكوز}\index{غلوكوز} أكثر بكثير --- وهو
        سكر \omterm{prop:g4:journey-of-food:absorb}{الدم} المغذي الجاهز، ووقودها اليومي؛
  \item وتطلق بالمقابل من \emph{\omterm{def:g7:what-is-respiration:respiration}{ثاني أكسيد الكربون}} أكثر ---
        ومن \emph{الحرارة}.
\end{itemize}
وهذا هو الاحتراق الخلوي في
\cref{prop:g7:what-is-respiration:power}، مقروءًا من عدّادات عضو:
وقود \omterm{def:g7:what-is-respiration:respiration}{وأكسجين} إلى الداخل، وعمل، وغاز فضلة ودفء إلى الخارج.
\end{proposition}
\admitted

\begin{example}[القياس عبر عضلة]\label{ex:g7:organs-at-work:measure}
يقارن علماء وظائف الأعضاء \omterm{prop:g4:journey-of-food:absorb}{الدم} \emph{الداخل} إلى \omterm{def:g3:bones-and-muscles:muscle}{عضلة} \omterm{prop:g4:journey-of-food:absorb}{بالدم}
\emph{الخارج} منها. ففي الراحة يكون \omterm{prop:g4:journey-of-food:absorb}{الدم} الخارج قد تخلى عن حصة
متواضعة من أكسجينه وغلوكوزه. وفي العمل الشاق تُظهر المقارنة نفسها
الدمَ الخارج مجرَّدًا --- وقد أُخذ أكثر أكسجينه القابل للتسليم ---
بينما قفز \omterm{def:g7:what-is-respiration:respiration}{ثاني أكسيد الكربون}. العضو نفسه والعدّادات نفسها: ولم
يتغير إلا العمل، وتبعته الفاتورة.
\end{example}

\begin{remark}[الدفء هو الإيصال]\label{rem:g7:organs-at-work:warmth}
حرارة الجسم العامل ليست عرضًا: فحرق الوقود يعطي دفئًا دائمًا،
والعضل غرفة أتون الجسم. والارتعاش هو الحيلة تُجرى عن قصد --- \omterm{def:g3:bones-and-muscles:muscle}{عضلات}
تُشغَّل لا لشيء إلا الدفء (وهي نفقة
\cref{ex:g3:balanced-diet:rest} الثابتة، مضاعفةً لحظةً). أما \omterm{def:g1:the-five-senses:sense}{جلد}
العدو المحمرّ فهو نظام التبريد يحمل حرارة الأتون بعيدًا --- وهو شأن
القسم التالي.
\end{remark}

\section{التوصيلات عند الطلب}

\begin{proposition}[الجسم يعيد تنظيم تموينه]\label{prop:g7:organs-at-work:supply}
لحظة بدء العمل تعيد خطوط التموين تنظيم نفسها لدفع الفاتورة
(وقد أعلنت ذلك \cref{ex:g4:heart-and-blood:effort}؛ وهذا هو
البلاغ الكامل):
\begin{enumerate}
  \item \emph{\omterm{def:g7:what-is-respiration:respiration}{الحركات التنفسية}} تعمق وتسرع --- فتعيد تحميل \omterm{prop:g4:journey-of-food:absorb}{الدم}
        أسرع عند \omterm{def:g7:how-animals-breathe:four}{الرئتين}
        (\cref{ex:g4:breathing:running})؛
  \item و\emph{\omterm{def:g4:heart-and-blood:heart}{القلب}} يدقّ أسرع وأقوى --- جولات أكثر في الدقيقة
        (\cref{prop:g4:heart-and-blood:pulse})؛
  \item \omterm{prop:g4:journey-of-food:absorb}{والدم} \emph{يُعاد توجيهه}: فتتسع \omterm{def:g4:heart-and-blood:vessels}{الأوعية} التي تخدم
        \omterm{def:g3:bones-and-muscles:muscle}{العضلات} العاملة (\omterm{def:g1:the-five-senses:sense}{والجلد} المبرّد)، بينما تضيق \omterm{def:g4:heart-and-blood:vessels}{أوعية}
        \omterm{def:g4:journey-of-food:tract}{الأنبوب الهضمي} المستريح --- \omterm{prop:g4:journey-of-food:absorb}{الدم} نفسه، محوَّلًا إلى الطابق
        المنفِق.
\end{enumerate}
وتتبع الثلاثة كلها العملَ تلقائيًّا، يوجّهها \omterm{def:g5:brain-in-command:system}{الدماغ} بخدماته التي لا
تُطلب (\cref{rem:g5:brain-in-command:more}).
\end{proposition}
\admitted

\begin{example}[البلاغ محسوسًا من الداخل]\label{ex:g7:organs-at-work:felt}
تفقّد حال عدو الدرج، مشروحًا سطرًا سطرًا: نفَس متسارع --- إعادة
تحميل؛ \omterm{def:g4:heart-and-blood:heart}{وقلب} يدقّ --- جولات أسرع؛ \omterm{def:g1:my-body:parts}{ورجلان} تحترقان --- فاتورة \omterm{def:g3:bones-and-muscles:muscle}{العضلات}
تُدفع عند الحدّ؛ \omterm{def:g1:the-five-senses:sense}{وجلد} دافئ محمرّ --- حرارة الأتون محمولةً إلى سطح
التبريد. بل حتى الغرزة الكلاسيكية والكسل بعد وجبة كبيرة يوافقان
ذلك: \omterm{def:g4:journey-of-food:digestion}{فالهضم} والعدو يتنافسان على \omterm{prop:g4:journey-of-food:absorb}{الدم} المحوَّل، ولهذا نصحت
\cref{met:g4:journey-of-food:habits} بألّا يُسابَق على \omterm{def:g4:journey-of-food:tract}{معدة} ملأى.
\end{example}

\begin{omfigure}
\begin{tikzpicture}
\begin{axis}[omaxis,
    width=9.6cm, height=5.8cm,
    xlabel={الدقائق}, ylabel={دقات القلب في الدقيقة},
    xmin=0, xmax=13.6, ymin=40, ymax=200,
    xtick={0,2,4,6,8,10,12},
    ytick={60,100,140,180},
    ]
  \addplot[very thick, omThm, smooth] coordinates {
    (0,72) (1,74) (2,120) (3,150) (4,158) (5,160)
    (6,120) (7,100) (8,88) (9,80) (10,76) (11,74) (12,73)
  };
  \draw[dashed, omExo] (axis cs:2,40) -- (axis cs:2,180);
  \draw[dashed, omExo] (axis cs:5,40) -- (axis cs:5,180);
  \node[font=\footnotesize] at (axis cs:1.0,52) {راحة};
  \node[font=\footnotesize] at (axis cs:3.5,52) {جهد};
  \node[font=\footnotesize] at (axis cs:8.6,52) {تعافٍ};
\end{axis}
\end{tikzpicture}

{\small \omterm{prop:g4:heart-and-blood:pulse}{نبض} عبر ثلاث دقائق من الجهد الشاق: الصعود مع العمل،
والاستواء قرب الأقصى، والعودة البطيئة --- \omterm{def:g7:organs-at-work:recovery}{التعافي} --- وقد صار
الجسم يسوّي حساباته.}
\end{omfigure}

\begin{definition}[التعافي]\label{def:g7:organs-at-work:recovery}
حين يتوقف الجهد لا يتوقف \omterm{prop:g4:heart-and-blood:pulse}{النبض} \omterm{def:g7:what-is-respiration:respiration}{والتنفس} معه: بل يعودان إلى ما كانا
عليه بالتدريج فحسب --- وهي مرحلة \emph{التعافي}\index{التعافي}.
والجسم يسوّي ديونًا: يعيد تحميل مخازن أكسجينه، وينظّف فضلات الجهد،
ويطرح الحرارة المتراكمة. وكلما زادت لياقة الجسم أسرع تعافيه ---
ولهذا كان زمن التعافي نفسه قياسًا للياقة.
\end{definition}

\section{التمرين، وحدود الآلة}

\begin{proposition}[التمرين يرفع السقف]\label{prop:g7:organs-at-work:training}
إذا مُورس التمرين بانتظام قويت سلسلة التموين كلها
(\cref{rem:g4:heart-and-blood:rest}،
\cref{met:g4:breathing:care}): \omterm{def:g4:heart-and-blood:heart}{فالقلب} يحرّك دمًا أكثر في الدقّة،
\omterm{def:g3:bones-and-muscles:muscle}{وعضلات} \omterm{def:g7:what-is-respiration:respiration}{التنفس} تعمق، \omterm{def:g3:bones-and-muscles:muscle}{والعضلات} نفسها تكبر وتدّخر وقودًا أكثر،
وأوعيتها الدقيقة تتكاثر. والنتيجة: الجهد نفسه يكلّف الجسم المدرَّب
نبضًا أقل، ويرتفع سقفه --- أي أشقّ جهد يستطيع أن يدفع ثمنه.
فالآلة يعيد بناءَها استعمالُها هي.
\end{proposition}
\admitted

\begin{example}[عدّاءان وتلّة واحدة]\label{ex:g7:organs-at-work:runners}
غير المدرَّب والمدرَّب يجريان التلّة نفسها. غير المدرَّب: \omterm{prop:g4:heart-and-blood:pulse}{نبض} عند
$180$ في منتصفها، ولهاث، واضطرار إلى المشي --- فخطوط التموين لا
تستطيع دفع الفاتورة، والدَّين يستدعي التوقف. والمدرَّب: \omterm{prop:g4:heart-and-blood:pulse}{نبض} $150$
عند القمة، والكلام ممكن، \omterm{def:g7:organs-at-work:recovery}{والتعافي} في ثلاث دقائق. التلّة نفسها،
والفاتورة نفسها --- وسقفان مختلفان، بنتهما شهور من تلال كهذه
تُؤخذ برفق.
\end{example}

\begin{method}[قراءة آلتك أنت]\label{met:g7:organs-at-work:read}
تجربة ذاتية بساعة
(و\cref{met:g4:heart-and-blood:pulse} توفّر تقنية أخذ
\omterm{prop:g4:heart-and-blood:pulse}{النبض}):
\begin{enumerate}
  \item سجّل \omterm{prop:g4:heart-and-blood:pulse}{نبض} الراحة، جالسًا هادئًا؛
  \item وقم بثلاث دقائق من جهد ثابت --- صعودًا على درجة، أو نطًّا
        بحبل؛
  \item وسجّل \omterm{prop:g4:heart-and-blood:pulse}{النبض} في الحال، ثم كل دقيقة حتى يعود إلى الراحة:
        وهو زمن \omterm{def:g7:organs-at-work:recovery}{التعافي}؛
  \item وأعد ذلك أسبوعيًّا عبر فصل دراسي من التمرين المنتظم:
        وراقب \omterm{prop:g4:heart-and-blood:pulse}{نبض} الراحة \omterm{prop:g4:heart-and-blood:pulse}{ونبض} الجهد وزمن \omterm{def:g7:organs-at-work:recovery}{التعافي} وهي تنجرف كلها
        إلى أسفل --- وهي إعادة بناء
        \cref{prop:g7:organs-at-work:training}، في معطياتك أنت.
\end{enumerate}
\end{method}

\section{تمارين}

\begin{exercise}[$\star$]\label{exo:g7:organs-at-work:1}
عدّد ما تأخذ منه \omterm{def:g3:bones-and-muscles:muscle}{العضلة} العاملة أكثر، وما تطلق منه أكثر.
\end{exercise}

\begin{exercise}[$\star$]\label{exo:g7:organs-at-work:2}
ما \omterm{prop:g7:organs-at-work:bill}{الغلوكوز}، ومن أين تناله \omterm{def:g3:bones-and-muscles:muscle}{العضلة} العاملة؟
\end{exercise}

\begin{exercise}[$\star$]\label{exo:g7:organs-at-work:3}
صف مقارنة \omterm{prop:g4:journey-of-food:absorb}{الدم} الداخل \omterm{prop:g4:journey-of-food:absorb}{بالدم} الخارج، وما تُظهره في الراحة وفي
الجهد.
\end{exercise}

\begin{exercise}[$\star$]\label{exo:g7:organs-at-work:4}
سمِّ إعادات تنظيم التموين الثلاث في
\cref{prop:g7:organs-at-work:supply}.
\end{exercise}

\begin{exercise}[$\star$]\label{exo:g7:organs-at-work:5}
لماذا يدفّئك العمل الشاق --- وما الارتعاش بمصطلحات هذا الفصل؟
\end{exercise}

\begin{exercise}[$\star$]\label{exo:g7:organs-at-work:6}
ما \omterm{def:g7:organs-at-work:recovery}{التعافي}، وأي ديون يسوّيها الجسم في أثنائه؟
\end{exercise}

\begin{exercise}[$\star\star$]\label{exo:g7:organs-at-work:7}
اقرأ منحنى \omterm{prop:g4:heart-and-blood:pulse}{النبض}: قيمة الراحة، وقيمة الاستواء، وكم يستغرق \omterm{def:g7:organs-at-work:recovery}{التعافي}
تقريبًا. وما الذي يغيّره التمرين في كل واحدة؟
\end{exercise}

\begin{exercise}[$\star\star$]\label{exo:g7:organs-at-work:8}
اشرح كسل ما بعد الغداء وغرزة العدّاء بقاعدة إعادة التوجيه.
\end{exercise}

\begin{exercise}[$\star\star$]\label{exo:g7:organs-at-work:9}
لماذا كان \omterm{def:g1:the-five-senses:sense}{جلد} الجهد المحمرّ تغيّرًا في \emph{التوصيل} لا خللًا؟
ومشكلة من يحلّ؟
\end{exercise}

\begin{exercise}[$\star\star$]\label{exo:g7:organs-at-work:10}
عدّد ثلاث إعادات بناء يرفع بها التمرين السقف، والعلامة اليومية
لكل واحدة.
\end{exercise}

\begin{exercise}[$\star\star$]\label{exo:g7:organs-at-work:11}
يجري تلميذان \cref{met:g7:organs-at-work:read}: نبضا راحة $68$
و$86$؛ وتعافيان ثلاث وثماني دقائق. فأي الآلتين المدرَّبة، وبم
يشهد كل من العددين بالضبط؟
\end{exercise}

\begin{exercise}[$\star\star\star$]\label{exo:g7:organs-at-work:12}
اربط ثلاثة فصول في فقرة واحدة: فاتورة \omterm{def:g3:bones-and-muscles:muscle}{العضلة}
(\cref{prop:g7:organs-at-work:bill})، ومقصود \omterm{def:g7:what-is-respiration:respiration}{التنفس}
(\cref{prop:g7:what-is-respiration:power})، وخدمة توصيل \omterm{prop:g4:journey-of-food:absorb}{الدم}
(\cref{prop:g4:heart-and-blood:carries}) --- حسابًا واحدًا لعدو
واحد.
\end{exercise}

\section{مسألة: دراسة اللياقة في المدرسة}

\begin{problem}[{مسألة نهاية الأسبوع --- فصل دراسي من معطيات
النبض، مقروءًا قراءة عالم وظائف}]\label{pb:g7:organs-at-work:1}
يجري الصف \cref{met:g7:organs-at-work:read} كل اثنين فصلًا
دراسيًّا: \omterm{prop:g4:heart-and-blood:pulse}{نبض} الراحة، \omterm{prop:g4:heart-and-blood:pulse}{والنبض} بعد ثلاث دقائق من الصعود على درجة، وزمن
\omterm{def:g7:organs-at-work:recovery}{التعافي}. وسجلّان مجهولا الهوية: التلميذ أ --- الأسبوع الأول: راحة
$84$، وجهد $168$، وتعافٍ $9$ دقائق؛ والأسبوع الثاني عشر: راحة
$72$، وجهد $150$، وتعافٍ $5$ دقائق. والتلميذ ب --- الأسبوع الأول:
راحة $70$، وجهد $150$، وتعافٍ $4$ دقائق؛ والأسبوع الثاني عشر: راحة
$69$، وجهد $148$، وتعافٍ $4$ دقائق.

\medskip\noindent\textbf{القسم الأول --- داخل جلسة واحدة.}
\begin{enumerate}
  \item في أثناء الصعود على الدرجة، عدّد ما يحدث لسحب \omterm{def:g3:bones-and-muscles:muscle}{عضلة} الفخذ
        من \omterm{def:g7:what-is-respiration:respiration}{الأكسجين} \omterm{prop:g7:organs-at-work:bill}{والغلوكوز}، ولإخراجها \omterm{def:g7:what-is-respiration:respiration}{ثاني أكسيد الكربون}
        والحرارة.
  \item اشرح لماذا يصعد \omterm{def:g7:what-is-respiration:respiration}{التنفس} \omterm{prop:g4:heart-and-blood:pulse}{والنبض} \emph{معًا} في دقائق
        الجهد.
  \item يشعر تلميذ بالدفء والاحمرار عند الدقيقة الثالثة. أعطِ
        حساب الأتون والتبريد.
  \item ولماذا لا يهبط \omterm{prop:g4:heart-and-blood:pulse}{النبض} إلى الراحة لحظة توقف الصعود؟
\end{enumerate}

\medskip\noindent\textbf{القسم الثاني --- قراءة الفصل الدراسي.}
\begin{enumerate}[resume]
  \item صف تغيرات أ عبر الفصل، وسمِّ إعادات البناء في
        \cref{prop:g7:organs-at-work:training} وراء كل
        عدد.
  \item أعداد ب لم تكد تتحرك. أعطِ القراءتين الصادقتين --- وأي
        واقعة إضافية عن حياة ب خارج الصف تحسم بينهما.
  \item في الأسبوع الخامس يقيس تلميذ بعد غداء متعجَّل مباشرة
        فيحصل على أعداد غريبة. فأي قاعدة من
        \cref{prop:g7:organs-at-work:supply} تدخّلت، وماذا ينبغي
        أن يقول البروتوكول عن أوقات القياس؟
  \item ولماذا يجب أن يقيس كل تلميذ بالطريقة نفسها دائمًا ---
        الجهد نفسه والتوقيت نفسه؟ (صديق قديم بين الطرائق؛
        سمِّه.)
\end{enumerate}

\medskip\noindent\textbf{القسم الثالث --- رسالة عالم
الوظائف.}
\begin{enumerate}[resume]
  \item اكتب ملخص فصل أ في جملتين للوحة إعلانات الصف، مستعملًا
        \emph{السقف} و\emph{\omterm{def:g7:organs-at-work:recovery}{التعافي}}.
  \item يسأل وليّ أمر أنبض الراحة المنخفض مقلق. أجب بدرّاج
        \cref{exo:g4:heart-and-blood:11}.
  \item يريد الصف عددًا واحدًا يتابعه في الفصل القادم. حاجج
        لزمن \omterm{def:g7:organs-at-work:recovery}{التعافي} في مقابل \omterm{prop:g4:heart-and-blood:pulse}{نبض} الجهد، مستعملًا
        \cref{def:g7:organs-at-work:recovery}.
  \item اختم بعبرة الدراسة في جملة واحدة: ما الذي أعاد بناء آلة
        أ؟
\end{enumerate}
\end{problem}
